% titel: Blatt 1, A3: Summen in Python

% Teil I: Werkzeuge, kein Lösungsweg.

Python-Programme zum Summieren natürlichter Zahlen abändern, ein vorgebtes
Programm für $n=1,2,3,\dots$ auswerten und die Gauß-Formel der Summe
beweisen.

\begin{werkzeugkasten}[Python: range und sum]
\begin{itemize}
  \item \texttt{range(21)} liefert die Zahlen $0,1,\dots,20$
        ($\to$ ``bis 20''). Für ``bis 29'' schreibt man \texttt{range(30)}.
  \item \texttt{sum(iterable)} addiert alle Elemente,
        z. B. \texttt{sum(range(30))} $=0+1+\dots+29$.
  \item \texttt{range(n+1)} steht für die Zahlen $0,1,\dots,n$.
\end{itemize}
\end{werkzeugkasten}

\begin{werkzeugkasten}[Gauß-Formel durch Induktion]
Die Summe der ersten $n$ Zahlen:
\begin{equation}
  \sum_{k=0}^{n} k = \frac{n(n+1)}2 .
\end{equation}
Damit ist \texttt{sum(range(n+1))-n*(n+1)/2} für jedes $n$ gleich $0$ --
das ist die zu beobachtende ``Auffälligkeit''.
\end{werkzeugkasten}

\begin{achtung}[Darauf achten]
\begin{itemize}
  \item \texttt{range} ist ``halb-offen'': \texttt{range(30)} endet bei
        $29$, \texttt{range(21)} bei $20$.
  \item In Python Mehrfachsummen sauber in eine \texttt{sum}-Funktion mit
        Generatorausdruck schreiben.
  \item Die Formel $n(n+1)/2$ ist für alle $n\in\mathbb N$ zu beweisen --
        Induktionsanfang und Induktionsschritt fehlen dürfen nicht getrennt
        sein.
\end{itemize}
\end{achtung}
